\section{Introduction}
Vision--language models are entering materials and chemistry workflows as readers of rendered structures,
diffraction data and figures, and benchmarks of that capability now exist for crystal renders and XRD images
alike \citep{xcrys2506,xrdbench2605}. When such a model fails, the remedy depends on where the failure sits:
a perception failure is fixed at the input, in the render protocol or the representation, while a reasoning
failure is fixed in the model. Separating the two requires an independent read of what an image contains,
and for natural photographs that reference does not exist: no procedure recovers a photograph's
information content without invoking a model of some kind. Rendered scientific figures are different. A
ball-and-stick crystal render is an orthographic projection of a structure whose coordinates are known
exactly, under a camera set the experimenter chose, and inverting that projection, by triangulating atom
positions across views and applying a deterministic symmetry algorithm to the reconstruction, returns the
answer the images support with no model at any step. This paper builds that instrument and reports what it
attributes. A two-stage pipeline, which replaces that inversion with a second model, has no model-free row at
all.

\begin{figure*}[t]
\centering
\resizebox{0.92\textwidth}{!}{%
\begin{tikzpicture}[
  font=\small,
  x=8.2cm, y=1cm,
  free/.style={circle, draw=black!70, fill=teal!70!black, inner sep=0pt, minimum size=5.4pt},
  mod/.style={circle, draw=black!70, fill=blue!55!black, inner sep=0pt, minimum size=5.4pt},
  lab/.style={anchor=east, inner sep=1pt},
  sub/.style={anchor=east, inner sep=1pt, black!50, font=\footnotesize},
  val/.style={anchor=west, inner sep=0pt, font=\footnotesize},
  iv/.style={black!45, {|[width=4.5pt]}-{|[width=4.5pt]}},
  ivl/.style={font=\footnotesize, black!55}]

\def\rA{2.10}\def\rB{1.20}\def\rC{0.00}\def\rD{-0.90}
\draw[black!25, dashed] (0.1429,-1.95) -- (0.1429,2.55);
\node[anchor=south, black!45, font=\footnotesize] at (0.1429,2.56) {chance};

\foreach \y in {\rA,\rB,\rC,\rD}{\draw[black!10] (0,\y) -- (1,\y);}

\node[lab] at (-0.03,\rA) {$\Rzero$\ \ true positions};
\node[sub] at (-0.03,\rA-0.36) {symmetry algorithm};
\node[free] at (1,\rA) {}; \node[val] at (1.015,\rA) {1.0000};

\node[lab] at (-0.03,\rB) {$\Rone$\ \ render geometry};
\node[sub] at (-0.03,\rB-0.36) {inversion, then algorithm};
\node[free] at (1.0,\rB) {}; \node[val] at (1.015,\rB) {1.0000};
\node[anchor=east, font=\footnotesize\bfseries, teal!55!black] at (0.9840,\rB) {no model};

\node[lab] at (-0.03,\rC) {$\Rthree$\ \ exact geometry as text};
\node[sub] at (-0.03,\rC-0.36) {model};
\node[mod] at (0.8524,\rC) {}; \node[val] at (1.015,\rC) {0.8524};

\node[lab] at (-0.03,\rD) {$\Rfour$\ \ five renders (pixels)};
\node[sub] at (-0.03,\rD-0.36) {model};
\node[mod] at (0.7333,\rD) {}; \node[val] at (1.015,\rD) {0.7333};

\draw[iv] (0.8524,0.62) -- (1.0,0.62);
\node[ivl, anchor=south] at (0.9262,0.70) {residual 0.1476};
\draw[iv] (0.7333,-1.42) -- (0.8524,-1.42);
\node[ivl, anchor=north] at (0.7929,-1.50) {perception 0.1191};

\draw[black!35] (0,-1.95) -- (1,-1.95);
\foreach \x/\t in {0/0, 0.25/0.25, 0.5/0.50, 0.75/0.75, 1/1.00}{
  \draw[black!35] (\x,-1.95) -- (\x,-2.07);
  \node[anchor=north, black!55, font=\footnotesize] at (\x,-2.09) {\t};}
\node[anchor=north, black!55, font=\footnotesize] at (0.5,-2.44) {crystal-system accuracy};
\end{tikzpicture}}
\caption{The attribution ladder ($n = 210$, $K{=}3$). Each rung answers the same question on the same
structures with one capability removed, and $\Rone$ answers it from the renders' own projection geometry with
no model in the measurement, which is what fixes the top of the ladder. With the extraction tolerance matched
to the symmetry tolerance, $\Rone$ and $\Rzero$ coincide at 1.0000, so the whole 0.2667 the strongest model
leaves is the model's. Model rungs are \rlLadder{} (Section~\ref{sec:roster}).}
\label{fig:teaser}
\end{figure*}

The shape of the argument has prior art. Several studies attribute multimodal failure to perception rather
than inference, whether by converting the image into a model-written description \citep{arc2512} or by
decoupling the two stages in post-training \citep{decouple2605}, but in every prior design we found the
perception stage is itself a model, so the split compares two model configurations, inherits the describer's
errors, and states no quantity for how much of the task the image answers at all. Our separation differs from
those closest in kind on three axes at once: the perception stage is a closed-form inversion rather than a
model or a written description; the reference is an accuracy on the same labels and structures, so it enters
the ladder as a rung; and the split is task-level rather than mechanistic, where attention-masking analyses
locate the layer at which visual access stops mattering \citep{vab2607}. What that buys is a number. Benchmark
headroom is otherwise taken to be one minus the best score, which assumes a ceiling of 1
\citep{zerobench2502}, and assumed ceilings can be wrong by enough to change conclusions
\citep{ceilingart2605}. Where a ceiling is measured rather than assumed, existing practice still puts a model
in it, typically a best-case ensemble over several models' predictions or a hypothetical selector computed
against the labels: an analytical device for what current models can reach rather than a measurement of what
the images contain. The render ceiling is the first such quantity computable with no model anywhere in the
measurement, and it comes with a characterisation of exactly when it is exact (Figure~\ref{fig:teaser}). Here we construct that ceiling, certify the condition under which it is exact, and use it as the top rung of an attribution ladder over fourteen vision-language models. Four contributions follow.

\textbf{(i) Perception is not the bottleneck for most of the roster, once the ceiling is measured.} Used as
the top rung of an attribution ladder on 210 structures under a frozen evaluation protocol, the oracle contradicts
the perception-bottleneck reading it was built to test: with exact geometry supplied as text, the residual
that survives is the larger deficit for 13 of 14 scored models. The single exception ranks second on pixels
and joint-first with exact geometry, so perception dominance is confined to the top of the roster rather than
being a property of the task, which reconciles this result with model-in-the-loop studies that examine
frontier models. We test for a monotone trend across the roster, do not find one at $n=14$, and report the
negative.

\textbf{(ii) A failure a leaderboard cannot see.} A strong model used as the extraction stage emits
syntactically perfect coordinate lists whose median recall against ground truth is 0.0000, a fabrication that
downstream accuracy alone would have scored as bad reasoning and that is detectable only against a model-free
reference.

\textbf{(iii) A characterisation of when a render determines its label, and of what the resulting ceiling
bounds.} Our identifiability result shows the oracle recovers exactly the atoms plus a set of geometric
phantoms, so it never drops an atom and its only failure mode is coincidence. That condition is met here, and
emptiness is a property of the released protocol rather than of the construction, since the same census finds
phantoms once cameras are removed. A companion scope result fixes the reading: the ceiling bounds an
exact-extraction centroid-triangulating pipeline and nothing outside that class. Because $\Rzero$ and $\Rone$
coincide at 1.0000, every point of a model's deficit is attributable to the model.

\textbf{(iv) What a practitioner does with the ceiling: a frontier and a conditioning constant.}
Identifiability being exact, what a practitioner faces is set by the tolerance $\delta$ at which their
extraction stage merges triangulated candidates, and the frontier we trace across $\delta$ says what an
imperfect reader of the same images would lose, which is narrower than what the images withhold. Camera
placement acts through the conditioning constant $\kappa$: fixed in advance for two protocols differing in
cameras alone, it predicted which would be less robust to extraction noise, and we report that prediction with two
that failed. Labels are computed rather than annotated, under a composition-exclusion protocol.

A sceptical reading is that an assumed ceiling of 1 was right all along. Measuring converts that assumption
into a certificate, the certificate fails at two and three views so emptiness is designable rather than given,
and at the merge tolerance our released pipeline used the same oracle reads $\RoneRel = 0.9524$, so an
instrument built on that pipeline alone would have charged 4.8 points of every model's deficit to the
images.

\subsection{Related work}
\label{sec:related}

\textbf{Model-in-the-loop attributions of multimodal failure.} The claim that multimodal failures are
perceptual rather than inferential is argued in several places, always by putting a second model in the loop,
whether as the describer, the decoupled perception stage, or the masked probe of the attribution studies
cited in the introduction. \citet{textoracle2512} serialise the visual stream into
text and conclude that the ceiling is set by the backbone's reasoning; their oracle is a ground-truth
\emph{description}, so it fixes the model's input but leaves the image's information content unmeasured.
\citet{mathlens2510} decompose performance into perception, reasoning and integration on geometry problems
whose diagrams and descriptions come from a common symbolic specification, and find integration dominant;
their perception measure is bundled together with integration rather than isolated. Additional reasoning can
even \emph{reduce} accuracy \citep{dual2509, cotdeg2604}. \citet{vdlm2404} convert raster vector graphics to
SVG with a rule-based encoder before a learned abstraction step; their symbolic stage is still learned, and
none of these designs reports a quantity for how much of the task the image answers on its own.

\textbf{Crystallography with vision--language models.} \citet{xcrys2506} stress-test nine VLMs on rendered
crystal structures under spatial- and com\-po\-si\-tion-exclusion protocols and report that supplying the source
crystallo\-graphic file alongside the image narrows but does not close the gap; the ceiling their gap is read
against is the best observed score rather than a measured bound. \citet{xrdbench2605} benchmark VLMs on XRD
peak indexing, establish that models fail on rendered crystallographic images, and likewise report that
supplying the source file alongside the image fails to close the gap on that target. On the model-free side,
\citet{cryspnet2003} and \citet{extinction2411} predict symmetry from composition and from
diffraction-consistent descriptors respectively, taking chemical or diffraction input rather than an image.
\citet{ziletti2018insightful} classifies crystal symmetry from computed diffraction images with a learned
network, so the reader itself is a trained model rather than a closed-form procedure.

\textbf{Resolution and render protocol.} \citet{resbench2510} and \citet{nativevis2506} show that dynamic
resolution input is a live failure mode for vision-language models. Multi-view correspondence is documented
as a weakness \citep{allangles2504}, and view selection has been optimised through a differentiable renderer
\citep{mvtn2011}, in both cases treating the render or view configuration as something to be learned or
searched rather than analysed for the conditioning it induces.

\textbf{Identifiability from multiple views.} Recovering shape from several orthographic views is classical
\citep{tomasi1992factorization}, and recent work does it with learned self-supervised objectives
\citep{gaussiancad2026}. On the representation-learning side, \citet{daunhawer2023identifiability} give
identifiability guarantees for the latent factors a multi-view contrastive objective recovers; the guarantee
concerns what a learned model recovers rather than exact geometric identifiability of scene content under a
known renderer.

\textbf{Verifying intermediate steps.} \citet{groundscore2603} verify explicit visual premises to make process
reward models reliable, and \citet{cvt2511} add continuous visual tokens for intermediate reasoning; both
treat the intermediate representation as something to be made checkable, without an exact ground-truth
reference against which to check it.

Across these lines, the shared gap is a model-free reference: every design above puts a second model, a
learned encoder, or a searched configuration where a closed-form inversion could stand instead, and none
reports a quantity for how much of the task an image answers on its own. Where the present findings and the
model-in-the-loop attributions overlap, the disagreement is one of measurement rather than finding: exact
geometry closes only a minority of the deficit for 13 of 14 models on the roster used here, and the two
readings are compatible once roster position is conditioned on, since the model-in-the-loop studies above
examine frontier models and the present perception-dominant exception is itself one.
